\section{横截性}
\label{sec:hamiltonian-floer-transversality}

一般而言，通过通有地选取哈密顿量 $H$ 或殆复结构 $J$，总能使\emph{光滑}\personFloer{}模空间达到横截。
但对紧化\personFloer{}模空间却不能如此，因为会出现负\personChern{}数多重覆叠球面的冒泡。一般
必须在\termKuranishiStructure{}的抽象框架中，利用虚基本链与多值扰动技术取得定义各类\personFloer{}映射所需的一般横截性。

\stfhIndex{I112}{Huofeier-Salameng}{霍费尔--萨拉蒙（Hofer--Salamon）}
我们把讨论限制在半正情形，依照\personHofer{}与\personSalamon{} (HS95) 的叙述，说明半正情形下
如何满足定义 19.2.7 和 19.3.2 中的假设。借助\termKuranishiStructure{}处理一般情形留给
(FOn99)。先给出 (HS95)、(On95) 中半正辛流形的定义。

\begin{definition}
\label{def:semipositive-symplectic-manifold}
若不存在 $A\in\pi_2(M)$ 使得
\[
 \omega(A)>0\qquad\text{且}\qquad 3-n\leq c_1(A)<0,
\]
则称辛流形 $(M,\omega)$ 是\emph{半正}的。
\end{definition}

对边界映射 $\partial_{(H,J)}$，有如下定理。

\begin{theorem}
\label{thm:semipositive-boundary-regularity}
假设不存在 $A\in\pi_2(M)$ 使得
\[
 \omega(A)>0\qquad\text{且}\qquad 3-n\leq c_1(A)<0.
\]
考虑非退化哈密顿量 $H:S^1\times M\to\mathbb R$。则对 $H$-正则的 $J$，
定义 19.2.7 中的假设成立。
\end{theorem}

\begin{proof}
边界映射 $\partial_{(H,J)}$ 的矩阵元由相对\personConley{}--\personZehnder{}指标满足
\begin{equation}
 \mu_H([z^-,w^-],[z^+,w^+])=1
\label{eq:19-index-one-boundary-moduli}
\end{equation}
的模空间
$\widetilde{\mathcal M}(H,J;[z^-,w^-],[z^+,w^+])$ 决定。我们已经证明，该空间是
有限多个 $C\in\pi_2(z^-,z^+)$ 所对应的
$\widetilde{\mathcal M}(H,J;z^-,z^+;C)$ 的并。

目标是证明
$\mathfrak j_\omega=C^\infty([0,1],\mathcal J_\omega)$ 中存在稠密的
$H$-正则 $J$ 子集。先通有地选取 $J$，使
$\widetilde{\mathcal M}(H,J;z^-,z^+;C)$ 在
$\mu_{(H,J)}(C)\leq0$ 时为空，在 $\mu_{(H,J)}(C)>0$ 时具有预期维数。这来自通有横截
性，其证明与第 10.4 节相同。
\footnote{译者注：原书这里笼统声称指标不大于零时模空间为空，但定义 19.2.7(1)
明确保留 $z_0=z_1,C=0$ 的定常指标零解。此句若不排除该情形便与定义冲突；译文
照录，并在此指出缺失的例外。}

\begin{exercise}
\label{ex:prove-hamiltonian-generic-transversality}
证明此横截性。
\end{exercise}

对这样的通有 $J$，$\widetilde{\mathcal M}(H,J;z^-,z^+;C)$ 是光滑一维流形。
指标条件~\eqref{eq:19-index-one-boundary-moduli} 与指标公式 (18.4.25)
\[
 \mu_H([z,w\#A])=\mu_H([z,w])-2c_1(A)
\]
蕴含 $z^-\neq z^+$，所以对每个这样的 $C$，$\mathbb R$ 在
$\widetilde{\mathcal M}(H,J;z^-,z^+;C)$ 上自由作用。故商
\[
 \mathcal M(H,J;z^-,z^+;C)
 =\widetilde{\mathcal M}(H,J;z^-,z^+;C)/\mathbb R
\]
是零维流形。矩阵元
\[
 n_{(H,J)}([z^-,w^-],[z^+,w^+])
\]
应当是这些 $C$ 所对应的基数
\[
 \#\mathcal M(H,J;z^-,z^+;C)
\]
之和。为了定义这个数，需要证明
$\widetilde{\mathcal M}(H,J;z^-,z^+;C)/\mathbb R$ 紧致。

\begin{proposition}
\label{prop:index-one-semipositive-compactness}
假设 $\mu_H([z^-,w^-],[z^+,w^+])=1$。则对通有的 $J$，
$\mathcal M(H,J;[z^-,w^-],[z^+,w^+])$ 紧致。
\end{proposition}

\begin{proof}
由\personGromov{}--\personFloer{}紧性，
$\widetilde{\mathcal M}(H,J;z^-,z^+;C)/\mathbb R$ 中的任意发散子列，一般有子列收敛到
\[
 u_\infty=u_1\#\cdots\#u_N;\qquad
 u_1(-\infty)=z^-,\quad u_N(\infty)=z^+.
\]
这是由尖点\personFloer{}流线 $u_\ell$ 组成的连通链，满足
\[
 u_\ell^0(-\infty)=z_\ell^-,\qquad
 u_\ell^0(\infty)=z_\ell^+;\qquad
 z_\ell^+=z_{\ell+1}^-,
\]
其中 $1\leq\ell\leq N$，并且每个 $u_\ell$ 形如
\begin{equation}
 u_\ell=u_\ell^0+\sum_{k=1}^{L_\ell}v_\ell^k.
\label{eq:19-cusp-floer-components}
\end{equation}
这里 $\{v_\ell^k\}_k$ 是有限个球面气泡，记
$C_\ell=[u_\ell]\in\pi_2(z_\ell^-,z_\ell^+)$。
\footnote{译者注：原书式 (19.6.21) 对每个分支都以上限 $L$ 计数气泡，但后文立即
改用逐分支的 $L_j$ 并令 $L=\sum_jL_j$；译文统一把本式上限写作 $L_\ell$。}

为证明 $\mathcal M(H,J;z^-,z^+;C)$ 紧致，我们要证明：对通有的 $J$，当
$\mu_{(H,J)}(C)=1$ 时，对任意 $(z^-,z^+)$ 与
$C\in\pi_2(z^-,z^+)$，都有 $N=1$ 且 $\{v_\ell^k\}=\varnothing$。

回忆，每个 $v_\ell^k$ 是在某点 $(\tau_k,t_k)$ 附着于 $u_\ell$ 的
$J_{t_k}$-全纯球面，其中先验地也允许 $\tau_k=\infty$。先考察 $\tau_k=\infty$，
即气泡附着在周期轨道像集上的情形。

考虑参数化模空间
\[
 \mathcal M^{\mathrm{para}}(J;\alpha)
 =\bigcup_{t\in[0,1]}\{t\}\times\mathcal M(J_t;\alpha).
\]
它的虚维数为
\[
 2c_1(\alpha)+2n-6+1=2c_1(\alpha)+2n-5.
\]
对通有 $J$，$\mathcal M^{\mathrm{para}}(J;\alpha)$ 中在某处单射曲线组成的开子集是
维数 $2c_1(\alpha)+2n-5$ 的光滑流形。因此，只要该模空间非空，就必有
\begin{equation}
 c_1(\alpha)\geq3-n.
\label{eq:19-simple-sphere-chern-lower-bound}
\end{equation}
现在假设 $J$ 对所有在某处单射曲线都满足此性质。

另一方面，对多重覆叠气泡，特别是\personChern{}数满足
$3-n\leq c_1(\alpha)<0$ 的气泡，一般不能通过扰动 $J$ 获得所需横截性。\emph{这里半正性
起关键作用}。在半正性假设下，式~\eqref{eq:19-simple-sphere-chern-lower-bound} 自动蕴含
$c_1(\alpha)\geq0$。若通有地选取 $J$，使该式对所有在某处单射曲线成立，那么由结构
定理 8.7.3，任何不在某处单射的曲线 $u$ 都是多重覆叠：
$u=\widetilde u\circ\varphi$，其中 $\varphi:S^2\to S^2$ 是次数 $d>1$ 的分歧覆叠，
$\widetilde u$ 在某处单射。因此
\[
 c_1(u)=d\,c_1(\widetilde u)\geq c_1(\widetilde u)
\]
（等号仅在 $c_1(\widetilde u)=0$ 时成立），从而
\[
 \operatorname{vir.dim}\mathcal M(J;[u])
 \geq\operatorname{vir.dim}\mathcal M(J;[\widetilde u]).
\]
显然 $\widetilde u$ 与 $u$ 的像集相同，所以若 $\widetilde u$ 不与
$\operatorname{Per}(H)$ 相交，$u$ 的像也不与之相交。这里略用记号，
$\operatorname{Per}(H)\subset M$ 表示 $H$ 的周期轨道的像集。因此，只需关注在某处
单射的气泡。

现在考虑评价映射
\[
 \operatorname{ev}_0:\mathcal M_1^{\mathrm{para}}(J;\alpha)
 \longrightarrow M,\qquad \operatorname{ev}_0(v,z)=v(z).
\]
$\mathcal M_1^{\mathrm{para}}(J;\alpha)$ 的虚维数比
$\mathcal M^{\mathrm{para}}(J;\alpha)$ 大 $2$，故为
$2c_1(\alpha)+2n-3$。必要时进一步扰动 $J$ 后，评价映射在在某处单射曲线的开子集上
与给定的一维子流形 $\operatorname{Per}(H)\subset M$ 横截。因此若
$3-2c_1(\alpha)>1$，即若
\begin{equation}
 c_1(\alpha)<1,
\label{eq:19-evaluation-avoids-periodic-orbits}
\end{equation}
则 $\operatorname{ev}_0$ 的像不与 $\operatorname{Per}(H)$ 相交。更多气泡更难出现的
理由留给读者说明。

结合~\eqref{eq:19-simple-sphere-chern-lower-bound} 与
\eqref{eq:19-evaluation-avoids-periodic-orbits} 可知，若
$3-n\leq c_1(\alpha)\leq0$，则 $X_H$ 的任意周期轨道的像上都不会发生冒泡。

由于每个 $u_k$ 的指标应至少为 $1$，有 $\mu_{(H,J)}([u_k])\geq1$。既然不能在
$(\pm\infty,*)$ 处冒泡，冒泡只能发生在某点
$(\tau_\ell^k,t_\ell^k)\in\mathbb R\times\mathbb R/\mathbb Z$。于是考虑评价映射
\[
 \operatorname{ev}_C:\mathcal M_L(H,J;z^-,z^+;C)\longrightarrow M^L,
\]
其中 $L$ 是标记点个数，$\mathcal M_L(H,J;z^-,z^+;C)$ 是带 $L$ 个标记点的\personFloer{}流线模空间。更确切地，
\[
 \mathcal M_L(H,J;z^-,z^+;C)
 =\widetilde{\mathcal M}(H,J;z^-,z^+;C)
 \times(\mathbb R\times S^1)^L/\sim,
\]
其中等价关系 $\sim$ 定义为
\footnote{译者注：原书定义等价关系时把标记点列写成
$(z_1,\ldots,z_L)$，下一行却立即以坐标 $(\tau_i,t_i)$ 使用它们；译文统一写出这些
标记点的柱面坐标。}
\[
 (u;((\tau_1,t_1),\ldots,(\tau_L,t_L)))
 \sim(u';((\tau'_1,t'_1),\ldots,(\tau'_L,t'_L)))
\]
当且仅当对某个 $\tau^0$，
\[
 u'(\tau,t)=u(\tau-\tau^0,t),\qquad
 (\tau_i',t_i')=(\tau_i+\tau^0,t_i).
\]

由 $\sim$ 的定义，$\operatorname{ev}_C$ 良定义。式~\eqref{eq:19-cusp-floer-components}
中的每个位形都是纤维积
\[
 \mathcal M_{L_j}(H,J;z_j,z_{j+1};C_j)_{\operatorname{ev}_{C_j}}
 \!\times_{(\operatorname{ev}_{\alpha_1},\ldots,
 \operatorname{ev}_{\alpha_{L_j}})}
 \bigl(\mathcal M_1^{\mathrm{para}}(\alpha_1^j),\ldots,
 \mathcal M_1^{\mathrm{para}}(\alpha_{L_j}^j)\bigr)
\]
中的元素，并且
\footnote{译者注：原书此处写
$C=C^0\#\sum_{j=1}^N\sum_{k=1}^{L_j}\alpha_k^j$，但本证明未定义 $C^0$，而紧接着
的指标公式逐一使用 $C_j$。译文依极限的 $N$ 段分解显式写出
$C_1\#\cdots\#C_N$。}
\[
 C=C_1\#\cdots\#C_N\#
 \sum_{j=1}^{N}\sum_{k=1}^{L_j}\alpha_k^j.
\]
当 $N\geq1$ 时，这个纤维积的虚维数为
\begin{equation}
 \sum_{j=1}^{N}\bigl(\mu_H(C_j)+2L_j\bigr)-(N-1)
 +\sum_{j=1}^{N}\sum_{k=1}^{L_j}
 \bigl(2c_1(\alpha_k^j)+2n-3\bigr)-2nL,
\label{eq:19-bubbled-stratum-virtual-dimension}
\end{equation}
其中 $L_j\geq0$ 且 $L:=\sum_{j=1}^{N}L_j$。（若 $N=0$，则必有 $z^-=z^+$；
由 $J$ 的选取，不可能附着气泡。此时，对某个 $A\in\pi_2(M)$ 必有
\[
 [z^+,w^+]=[z^-,w^-\#A],\qquad C=C^0=-A.
\]
于是由指标公式，$\mu_H(z^-,z^+;C^0)=2c_1(A)$，不可能等于所要求的 $1$。）
此外，
\begin{equation}
 \mu_H(C)=\sum_{j=1}^{N}
 \left(\mu_H(C_j)+\sum_{k=1}^{L_j}2c_1(\alpha_k^j)\right).
\label{eq:19-total-index-bubble-decomposition}
\end{equation}

\begin{lemma}
\label{lem:index-one-bubbled-stratum-dimension}
假设 $\mu_H(C)=1$。则~\eqref{eq:19-bubbled-stratum-virtual-dimension} 中的虚维数为
$2-L-N$，其中 $N\geq1$。
\end{lemma}

\begin{proof}
把 $\mu_H(C)=1$ 代入~\eqref{eq:19-total-index-bubble-decomposition}，得到
\[
 \sum_{j=1}^{N}
 \left(\mu_H(C_j)+\sum_{k=1}^{L_j}2c_1(\alpha_k^j)\right)=1.
\]
再代入~\eqref{eq:19-bubbled-stratum-virtual-dimension}，并使用
$L=\sum_{j=1}^{N}L_j$，虚维数化为
\[
 1+2L-(N-1)+L(2n-3)-2nL=2-L-N.
\]
\end{proof}

若模空间横截，这个维数必须非负，故 $0\leq2-N-L$。因为 $N\geq1$，$L$ 只能为
$0$ 或 $1$。当 $L=0$ 时，$1=\sum_{j=1}^{N}\mu_H(C_j)$。又因
$\mu_H(C_j)\geq0$，除一个 $j$（不妨为 $j=1$）外，其余 $\mu_H(C_j)$ 都为 $0$。
于是对所有 $j\neq1$ 都有 $z_j^-=z_j^+$，且对 $j=2,\ldots,N$，至少有一个气泡
必须附着于 $z_j^-=z_j^+$；由 $J$ 的选取，这不可能。因此此时也有 $N=1$。

若 $L=1$，上述不等式给出 $N\leq1$，故必有 $N=1$。这意味着气泡附着于一个周期
轨道，但这已经被前面对 $J$ 的选取排除。因此必有 $N=1$、$L=0$；也就是说，
该序列既不产生非常值气泡，也不分裂为断裂流线。这就证明了当
$\mu_H(z^-,z^+;C)=1$ 时，每个 $\mathcal M(H,J;z^-,z^+;C)$ 都紧致。
\footnote{译者注：本段的 $L$ 是有限柱面点处的标记点（冒泡点）总数，所以
$N=1,L=1$ 描述的气泡一般附着于\personFloer{}流线内部，并不自动附着于周期轨道。
原书从 $N=1,L=1$ 直接推出“气泡附着于周期轨道”，与此前的标记点模空间设置
不符；正文照录该推断，但紧性证明在此存在缺口。}

由此命题，边界映射 $\partial_{(H,J)}$ 的矩阵元良定义，因而边界映射本身良定义。
类似的论证证明，当 $\mu_H(z^-,z^+;C)=2$ 时，
$\mathcal M(H,J;z^-,z^+;C)$ 除了上述讨论中 $N=2,L=0$ 的分裂外是紧致的。这就
证明了定义 19.2.7 中的假设，从而定理 19.6.2 得证。
\end{proof}
\end{proof}

\begin{exercise}
\label{ex:index-two-boundary-compactification-details}
补全上述证明最后一段中关于 $\mu_H(z^-,z^+;C)=2$ 情形的论证细节。
\end{exercise}

下面考虑链映射 $h_{(\mathcal H,\mathfrak j;\rho)}$。唯一的差别是这里不取
$\mathbb R$ 的商，而直接考虑模空间
\[
 \mathcal M((\mathcal H,\mathfrak j;\rho);[z_0,w_0],[z_1,w_1])
 =\bigcup_C\mathcal M((\mathcal H,\mathfrak j;\rho);z_0,z_1;C).
\]
其维数为 $\mu_{(\mathcal H,\mathfrak j;\rho)}(C)$，而定义链映射所需的是
$\mu_{(\mathcal H,\mathfrak j;\rho)}(C)=0$ 的情形。类似地可以证明如下定理。

\begin{theorem}
\label{thm:semipositive-chain-map-regularity}
假设 $(M,\omega)$ 中不存在 $A\in\pi_2(M)$ 使得
\[
 \omega(A)>0\qquad\text{且}\qquad3-n\leq c_1(A)<0.
\]
则定义 19.3.2 中的假设成立。
\end{theorem}

\begin{proof}
由同类论证，对 $N\geq1$ 的极限位形得到虚维数
\[
 \sum_{j=1}^{N}(\mu_H(C_j)+2L_j)-(N-1)
 +\sum_{j=1}^{N}\sum_{k=1}^{L_j}
 (2c_1(\alpha_k^j)+2n-3)-2nL,
\]
以及
\[
 \mu_H(C)=\sum_{j=1}^{N}\mu_H(C_j)
 +\sum_{j=1}^{N}\sum_{k=1}^{L_j}2c_1(\alpha_k^j).
\]
若 $N=0$，则必有
\[
 [z^-,w^-]=[z^+,w^+],\qquad C^0=0,
\]
因为极限位形中不可能附着任何气泡。因此，相应模空间
$\mathcal M_{(\mathcal H,\mathfrak j;\rho)}([z^-,w^-],[z^+,w^+])$
\footnote{译者注：链映射连接的两端一般分别属于
$\operatorname{Crit}\mathcal A_{H^0}$ 与
$\operatorname{Crit}\mathcal A_{H^1}$，原书却在 $N=0$ 时直接写
$[z^-,w^-]=[z^+,w^+]$，未说明如何识别这两个不同临界集；正文照录。}
紧致。

现设 $N\geq1$。这次考察 $\mu_H(C)=0$，得到
\[
 0=1-N-L.
\]
这蕴含 $L=0$ 且 $N=1$，进而
$\mathcal M_{(\mathcal H,\mathfrak j;\rho)}([z^-,w^-],[z^+,w^+])$ 紧致。
\end{proof}

接着研究 $h_{(\mathcal H,\mathfrak j;\rho)}$ 的链映射性质，即恒等式
\begin{equation}
 h_{(\mathcal H,\mathfrak j;\rho)}\circ\partial_{(H^0,J^0)}
 =\partial_{(H^1,J^1)}\circ h_{(\mathcal H,\mathfrak j;\rho)}.
\label{eq:19-chain-map-property}
\end{equation}
为此，考虑使链映射模空间
$\mathcal M_{(\mathcal H,\mathfrak j;\rho)}([z^-,w^-],[z^+,w^+])$ 虚维数为 $1$
的对子 $[z^-,w^-],[z^+,w^+]$，并分析其紧化边界。

对链同伦映射，考虑参数化模空间
\[
 \mathcal M^{\mathrm{para}}
 (\overline{\mathcal H},\overline{\mathfrak j};\rho;
 [z^-,w^-],[z^+,w^+]),
\]
其中
$\overline{\mathcal H}=\{\mathcal H_\kappa\}_{0\leq\kappa\leq1}$ 与
$\overline{\mathfrak j}=\{\mathfrak j_\kappa\}_{0\leq\kappa\leq1}$ 是单参数族，
\[
 \mathcal H_\kappa=\{H^{s,\kappa}\},\qquad
 \mathfrak j_\kappa=\{J^{s,\kappa}\};\qquad
 H^{0,\kappa}=H^\alpha,\quad H^{1,\kappa}=H^\beta
\]
对所有 $0\leq\kappa\leq1$ 成立。考虑方程族
\[
 \frac{\partial u}{\partial\tau}
 +J_t^{\rho(\tau),\kappa}
 \left(\frac{\partial u}{\partial t}-X_{H_t^{\rho(\tau),\kappa}}\right)=0,
 \qquad0\leq\kappa\leq1.
\]
考察这样的对子 $([z^-,w^-],[z^+,w^+])$：相应参数化模空间的虚维数为 $0$ 或
$1$，即
\[
 \mu_{H^\alpha}([z^-,w^-])-\mu_{H^\beta}([z^+,w^+])+1=0
 \quad\text{或 }1.
\]
维数为 $0$ 时，
\[
 \mu_{H^\alpha}([z^-,w^-])-\mu_{H^\beta}([z^+,w^+])=-1.
\]

\begin{lemma}
\label{lem:parameterized-index-minus-one-compactness}
对通有族 $(\overline{\mathcal H},\overline{\mathfrak j})$，当
\[
 \mu_{H^\alpha}([z^-,w^-])-\mu_{H^\beta}([z^+,w^+])=-1
\]
时，$\mathcal M^{\mathrm{para}}
(\overline{\mathcal H},\overline{\mathfrak j};\rho;
[z^-,w^-],[z^+,w^+])$ 紧致。
\end{lemma}

\begin{exercise}
\label{ex:fredholm-setup-parameterized-compactness}
表述正确的\personFredholm{}设置并证明此引理。
\end{exercise}

借助与这一同伦相伴的族 $\overline{\mathfrak j}$ 与
$\overline{\mathcal H}$，通过计数满足上式指标条件的参数化模空间中的元素，构造链同伦
映射
\[
 \Upsilon_{(\overline{\mathfrak j},\overline{\mathcal H};\rho)}:
 CF_*(J^0,H^0)\longrightarrow CF_{*+1}(J^1,H^1).
\]
特别地，映射 $\Upsilon_{(\overline{\mathfrak j},\overline{\mathcal H};\rho)}$ 的次数为
$1$。
\footnote{译者注：原书同一行把映射的陪域写为 $CF_{*+1}$，且此前第 19.4 节也由
指标差 $-1$ 推出链同伦映射次数为 $1$；但紧接着却称其“次数为 $-1$”。译文按映射
类型和第 19.4 节改为 $1$。}

现在对满足
\[
 \mu_{H^\alpha}([z^-,w^-])-\mu_{H^\beta}([z^+,w^+])=0
\]
的对子考虑上述参数化模空间。此时它的维数为 $1$，其边界如命题 19.4.1 所述。

\begin{remark}
\label{rem:general-case-virtual-moduli-cycle}
对一般辛流形，需要在\termKuranishiStructure{}的框架中使用虚模闭链和抽象多值扰动的概念
(FOn99)。
\end{remark}
